% C.3bn. Full root and both peer proof/source audits passed; stated scope unchanged.
\paragraph{C.3bn. Endpoint lifting and the ordinary top quotient
in all resolved dimensions (Proved under the stated hypotheses).}
Let $k$ be an algebraically closed field of characteristic $p>0$.
Let $R$ be an excellent henselian local coefficient-$k$ algebra,
with residue field $k$, obtained by strict localization of a
finite-type $k$-scheme.  Let
\[
 \pi:Z\longrightarrow\operatorname{Spec}R
\]
be projective, with $Z$ regular of pure dimension $d+1$, $d\geq1$,
and with $\pi$ an isomorphism away from the closed point.  Suppose
that the reduced closed fibre $E$ is the full support of an effective
Cartier divisor $D$ and that $\mathcal O_D(-D)$ is ample.  No
normal-crossing or smoothness hypothesis on $E$ is imposed.
Write $X=Z\setminus E$.

For every smooth connected projective curve $P/k$, two points
$o,t\in P(k)$, and class
\[
             \beta\in F_\gamma^dK_0(D\times P)_{\mathbf Q},
\]
there is an actual class $b\in K_0(Z\times P)_{\mathbf Q}$ such that
\begin{equation}
       b|_{E\times\{u\}}=\beta|_{E\times\{u\}}
                    \qquad(u=o,t).                    \tag{AE.1}
\end{equation}
All $K_0$ groups here are perfect $K$-groups, with the vector-bundle
gamma filtration.  The equality is at the two specified endpoints;
the prescribed class on the entire exceptional product need not lift.

\emph{1. A perfect determinantal representative of a gamma class.}
Put $T=D\times P$, of pure dimension $d+1$.  This is a projective
Cohen--Macaulay $k$-scheme.  Indeed, $D$ is Cartier in the regular
$Z$; a power of $\mathfrak m$ annihilates $\mathcal O_D$, and
$R/\mathfrak m^M$ is finite-dimensional over $k$.  Thus projectivity
over that Artinian coefficient algebra gives projectivity over $k$.
The same applies to $E$.  In particular these schemes have ample
line bundles and their perfect $K_0$ is vector-bundle $K_0$.

All rank calculations may be made componentwise on the finitely
many connected components of $T$.
For a rank-$n$ bundle $V$ on $T$ and $1\leq i\leq n$, set
$m=n-i+1$ and take a homomorphism $\mathcal O_T^m\to V$.
Suppose that its maximal-minor ideal has grade $i$ wherever it is
proper.  The Eagon--Northcott complex resolves the structure sheaf
of its degeneracy locus $C_i$.  This exact noetherian-ring statement
is \citet[Theorem 2.16, pp.~18--20; construction in Section 2.C,
pp.~16--18]{D_BrunsVetter1988Determinantal}.
In the natural global complex the term preceding degree zero in
homological degree $q-m+1$ is
\[
 \bigwedge^q V^*\otimes
       \bigl(\operatorname{Sym}^{q-m}\mathcal O_T^m\bigr)^*,
                       \qquad m\leq q\leq n.
\]
The determinant of the trivial rank-$m$ bundle is suppressed here.
Dual symmetric powers, rather than an identification with symmetric
powers of the dual, are used in positive characteristic.  Locally
this is precisely the natural contraction complex of the cited
construction.  Consequently its perfect class is
\begin{equation}
 \epsilon_i(V):=[\mathcal O_{C_i}]
  =1+\sum_{q=m}^n(-1)^{q-m+1}
       {q-1\choose m-1}[\bigwedge^qV^*].               \tag{AE.2}
\end{equation}
If the degeneracy locus is empty the same complex is exact and
the displayed class is zero.

There is the following exact universal identity:
\begin{equation}
       [\det V]\cdot\epsilon_i(V)
               =\sum_{j=i}^n\gamma^j([V]-n).          \tag{AE.3}
\end{equation}
To verify it, use the injective flag-bundle splitting principle,
write the line roots as $\ell_a=1+x_a$, and expand elementary
symmetric functions.  Equivalently multiply (AE.2) by $\det V$
and replace $q$ by $n-h$.  The coefficient identity is
\[
 \sum_{j=h}^{i-1}(-1)^{j-h}{n-h\choose j-h}
       =(-1)^{i-1-h}{n-h-1\choose i-1-h}
                  \quad(0\leq h<i).
\]
It is the partial alternating binomial sum.  The terms of degree
less than $i$ in $\prod_a(1+x_a)$ are exactly subtracted, giving
(AE.3).  No assertion about injectivity on associated graded groups
is needed.  It follows directly in $K_0(T)$ that
\begin{equation}
 \epsilon_i(V)\in F_\gamma^i,
 \qquad \epsilon_i(V)\equiv\gamma^i([V]-n)
                             \pmod {F_\gamma^{i+1}}.   \tag{AE.4}
\end{equation}
This is a perfect-complex identity and retains any Cartan-kernel
class on the nonreduced scheme.

\emph{2. Generate the top gamma piece by perfect curves.}
Modulo $F_\gamma^{d+1}K_0(T)$, the group $F_\gamma^d$ is generated
by finite products
\begin{equation}
          \prod_{a=1}^r\gamma^{i_a}([V_a]-n_a),
       \qquad i_a\geq1,\quad \sum_a i_a=d,             \tag{AE.5}
\end{equation}
where the $V_a$ may be chosen with sufficiently generated first
jets on a fixed finite smooth stratification of $E\times P$ and
on the two endpoint stratifications.  Here is the reduction from
arbitrary virtual bundles, so that (AE.5) is not restricted to a
single gamma operation on an actual bundle.  If $L$ is very ample
and $V$ any vector bundle, choose $M$ with $V^*\otimes L^M$
generated.  Dualizing its generating sequence and twisting gives
\[
 0\longrightarrow V\longrightarrow (L^M)^N
        \longrightarrow Q\otimes L^M\longrightarrow0,
\]
where $Q$ is globally generated.  Products of generating sections
of $Q$ with first-jet generators of $L^M$ generate first jets of
the last bundle.  Taking $M$ larger and applying this to finitely
many bundles writes each rank-zero virtual bundle as $A-B$ with
$A,B$ equally ranked and sufficiently positive.  A common positive
summand can make their ranks at least $d$.  Expand
\[
       \gamma_z(A-B)=\gamma_z(A-\operatorname{rk}A)
                       \gamma_z(B-\operatorname{rk}B)^{-1}.
\]
The degree-$d$ coefficient is a finite integral sum of products of
the form (AE.5).  Products of total degree greater than $d$ may be
discarded modulo the next filtration.  This proves the reduction.

For each factor choose independently a general matrix of sections
\[
              \mathcal O_T^{n_a-i_a+1}\longrightarrow V_a.
\]
On a smooth reduced stratum the rank-drop-by-one condition has
codimension $i_a$; the next rank drop has codimension
$2(i_a+1)$.  The choices are simultaneously general on finitely
many strata, including the two endpoint fibres.  First-jet
generation gives transversality of every intersection of the
rank-drop-by-one loci.  To spell out the characteristic-free
incidence bound, on a smooth stratum of dimension $q$ a total
rank-drop condition of codimension $c\leq q$ has failure of
transversality in codimension $c+(q-c+1)=q+1$ in the jet space.
If $c>q$, its ordinary incidence already misses a general choice.
The images of the bad incidences have closures of dimension
strictly less than the parameter space.  Their finite union is
avoidable over the infinite field $k$.

The final intersection $C=\bigcap_aC_{i_a}$ avoids every deeper
rank locus.  Indeed for the $a$th matrix the required incidence
has codimension at least
\begin{equation}
             2(i_a+1)+\sum_{b\ne a}i_b
                    =d+i_a+2>d+1=\dim T.             \tag{AE.6}
\end{equation}
Near $C$, row and column operations thus replace each matrix by
an identity block and $i_a$ functions.  Together these are $d$
functions whose restrictions have independent differentials on
the smooth stratum through the point.  They cut the expected
codimension $d$ in the Cohen--Macaulay scheme $T$, hence form a
regular sequence.  Each separate maximal-minor ideal has height
at least its expected value by the same stratified dimension
argument.  The opposite inequality is the determinantal height
bound\citep[Theorem 2.1 and proof, pp.~10--11]
{D_BrunsVetter1988Determinantal}; Cohen--Macaulayness then makes
grade equal to that height.  Thus the Eagon--Northcott resolutions
used in (AE.2) apply globally.
Their derived tensor product resolves $\mathcal O_C$: near $C$
it is the Koszul resolution of the combined regular sequence;
away from $C$ at least one factor is exact.  Therefore
\[
                  \prod_a\epsilon_{i_a}(V_a)=[\mathcal O_C]
                         \quad\hbox{in perfect }K_0(T).
\]
Thus $\beta$, modulo $F_\gamma^{d+1}$, is a finite rational
linear combination of these perfect classes of curves.

The reduction of $E$, although possibly singular, admits a finite
smooth locally closed stratification over the perfect field $k$.
In $E\times P$ the curve $C$ meets only strata of dimension at
least $d$, and its $d$ defining differentials are independent on
each stratum that it meets.  At the endpoints it meets only the
smooth top-dimensional strata of $E\times\{u\}$, and
\begin{equation}
         C_u\subset D\quad(u=o,t)
\end{equation}
is finite and a regular immersion of codimension $d$.  It may
be nonreduced; this causes no problem below.  We do not require
$C\to P$ flat away from the two selected endpoints.

\emph{3. A hypersurface flag leaves an actual regular threefold.}
Fix such a curve $C\subset T$.  Its conormal bundle
$I_C/I_C^2$ has rank $d$.  Take $d-1$ sufficiently high line
bundles $\mathcal A_j$ on $Z\times P$, ample over
$\operatorname{Spec}R$, and sections of
$I_C\mathcal A_j|_T$.  The line bundles can be chosen as exterior
products of a relatively ample bundle on $Z$ and an ample bundle
on $P$.  Serre vanishing gives the surjections onto the twisted
conormal sections on $C$, and enough first jets away from $C$.
For a general choice, the resulting $d\times(d-1)$ conormal
matrix has full rank everywhere on $C$.  The rank-defective
condition has codimension two, larger than $\dim C=1$.
It suffices to verify this on $C_{\mathrm{red}}$; Nakayama's
lemma then gives a locally split injection on $C$ itself.

Make the choices also transverse to all smooth strata away from
$C$, and transverse on the two endpoint stratifications.  Each
prefix of the flag can be required to have the corresponding
transversality.  The same first-jet dimension argument as in
Step~2 supplies these finitely many open conditions.  Along $C$
the full-rank conormal matrix, composed with the independent
$d$ differentials already obtained, gives $d-1$ independent
ambient differentials; the same is true after fixing either
endpoint.  Relative Serre vanishing lifts all these sections to
actual sections of $\mathcal A_j$ on $Z\times P$, since
\[
 H^1\bigl(Z\times P,\mathcal A_j(-D\times P)\bigr)=0.
\]
Let $W$ be their simultaneous zero scheme, and put
\begin{equation}
              H=W\cap(D\times P).                    \tag{AE.7}
\end{equation}
When $d=1$ the flag is empty, so $W=Z\times P$ and $H=T$.

The source $W$ is regular of pure dimension three along its
closed fibre, by the independent ambient differentials.  It
is regular everywhere: the nonregular locus is closed by
excellence, and its proper image in the local scheme
$\operatorname{Spec}R$ would contain the closed point if
nonempty.  No component of a nonempty closed subset of this
proper scheme can miss that fibre.  This argument also excludes
components of an unintended dimension.  Thus
$W\hookrightarrow Z\times P$ is a regular immersion of
codimension $d-1$ everywhere.  The same flag gives the expected
dimension two for its intersection with the reduced closed
fibre.  Hence $D\times P$ contains no component of $W$ and its
restriction is a Cartier divisor $H$ on $W$.

The curve $C$ is Cartier in $H$.  Locally a choice of basis in
the regular-sequence conormal identifies the $d-1$ flag equations
with the first $d-1$ generators of $I_C$; the last generator
is a nonzerodivisor on $H$ and cuts out $C$.  This statement
uses the actual multiplicities of $D$, not just its reduction.
At $u=o,t$, the endpoint transversality makes $W_u$ regular of
dimension two, $H_u$ Cartier in $W_u$, and $C_u$ Cartier in
$H_u$.  The equations of the flag, of the parameter point, and
of $D$ give the requisite regular sequences in either order.
In particular every pullback square used below is
Tor-independent.  These conditions also give flatness of $W$
and $H$ over $P$ near the two selected fibres.  Global flatness
over all of $P$ is unnecessary; vertical fibres elsewhere are
allowed.

The first-jet constructions take place on the finite-type
projective schemes $D\times P$ and their smooth strata.  The
ambient regularity assertions may equally be checked after
descending the finite projective presentation and its finitely
many equations to one pointed etale finite-type model of $Z$.
After shrinking around its proper exceptional fibre, that
regular model is smooth over $k$.  Thus no smoothness assertion
for an unspecified non-finite-type ambient variety has entered
the Bertini argument.

\emph{4. Lift a finite tensor power on the threefold section.}
Put $J=\mathcal O_W(-H)$ and $L=J|_H$.
The line bundle $L$ is ample relative to $H\to P$, being the
restriction of $\mathcal O_D(-D)$ from $D\times P$.
Relative Serre vanishing gives
$R^i(H\to P)_*L^n=0$ for $i>0$ and all $n$ beyond one bound.
Since $P$ has cohomological dimension one for coherent sheaves,
the Leray sequence gives
\begin{equation}
                         H^2(H,L^n)=0\quad(n\gg0).    \tag{AE.8}
\end{equation}
This argument does not require all fibres of $H\to P$ to be
curves.  The precise projective Serre input is
\citet[Lemma 30.17.1]{frC3azStacksSerre}, applied on a finite
affine cover of $P$ and taking the maximum of the bounds.

Start with $Q_0=\mathcal O_H(-C)$.  On the thickenings
$H_n=V(J^n)$, the obstruction to extending a line bundle from
$H_n$ to $H_{n+1}$ belongs to $H^2(H,L^n)$, by the square-zero
unit sequence\citep[Lemma 37.4.1]{frC3azStacksPicThickening}.
Each such group is killed by $p$.  At the finitely many steps
before (AE.8) applies, replace the partial bundle and all its
preceding restrictions by their $p$th tensor powers and extend.
There results one finite integer $e$ and a compatible formal
line bundle restricting to $Q_0^{p^e}$ on $H$.

The powers of $J$ and $\mathfrak m\mathcal O_W$ are cofinal,
because their supports are exactly the closed fibre.  Proper
Grothendieck existence algebraizes the formal line bundle on
$W_{\widehat R}$; algebraizing its inverse and using full
faithfulness verifies invertibility
\citep[Lemma 30.24.3]{frC3azStacksExistence}.
Regularity of the completion map and Popescu express
$\widehat R$ as a filtered colimit of smooth $R$-algebras
\citep[Lemma 15.51.14; Theorem 16.12.1]
{frC3azStacksCompletion,frC3azStacksPopescu}.
The bundle descends to one such algebra
\citep[Lemma 32.10.3(2)]{frC3azStacksLimitBundles}.
Choose $M$ with $\mathfrak m^M\mathcal O_W\subset J$.
The specified point of this smooth algebra over
$\widehat R/\mathfrak m^M=R/\mathfrak m^M$ lifts to $R$
by henselian smooth lifting for the pair $(R,\mathfrak m^M)$
\citep[Lemma 15.13.3]{frC3azStacksHenselSmooth}.
Pulling back along that lift gives an actual line bundle $Q$
on the fixed $W$ satisfying
\begin{equation}
                            Q|_H\simeq Q_0^{p^e}.      \tag{AE.9}
\end{equation}
The specified finite-jet point, rather than merely a residue-field
point, is preserved in this descent.

\emph{5. Exact perfect identities at finite endpoints.}
Let $j:W\hookrightarrow Z\times P$ and set
$b_C=j_*(1-[Q])$.  This is a perfect $K_0$ class, since $j$ is
a regular immersion.  The regular sequences in Step~3 give
derived base change at each $u=o,t$ and then along $D$.
Writing $i_u:H_u\hookrightarrow D$, this yields
\[
       b_C|_{D\times\{u\}}
             =(i_u)_*\bigl(1-[Q_{0,u}^{p^e}]\bigr).
\]
For any positive integer $a$, the Cartier sequences on $H_u$
give the identity
\begin{equation}
 (i_u)_*\bigl(1-[\mathcal O_{H_u}(-aC_u)]\bigr)
       =\sum_{j=0}^{a-1}
            [\mathcal O_{C_u}\otimes\mathcal O_{H_u}(-jC_u)]
       =a[\mathcal O_{C_u}]\quad\hbox{in }K_0(D).       \tag{AE.10}
\end{equation}
The second equality is exact: every line bundle on the finite
Artinian $k$-scheme $C_u$ is trivial.  All terms are perfect
on $D$ because the immersions are regular.  This is not a
cycle-only identity, and no claim about gamma compatibility
of pushforward is used.  Taking $a=p^e$ proves that
$p^{-e}b_C$ has the same two restrictions to $D$ as
$[\mathcal O_C]$.

Combine the finitely many curves from Step~2.  The omitted terms
belong to $F_\gamma^{d+1}K_0(T)_{\mathbf Q}$, whose restriction
to $E$ is zero by the dimension theorem
\citep[Chapter II, Remark 8.8.5, p.~83]{D_Weibel2013KbookII}.
It applies to the possibly singular projective $d$-fold $E$.
This proves (AE.1).

\paragraph{The corresponding presentation of the ordinary top quotient
(proved consequence in the same resolved scope).}
Define
\[
 S_d(E)=F_\gamma^dK_0(E)_{\mathbf Q},
\]
and let $J_K^d(E)$ be the subgroup generated by endpoint differences
of classes in $F_\gamma^dK_0(E\times P)_{\mathbf Q}$, with $P$
smooth connected projective and with actual product sections.
Restriction and the rational Chern character give an isomorphism
\begin{equation}
         \rho_d:\mathrm{CH}^d(Z)_{\mathbf Q}
                         \xrightarrow{\sim}S_d(E).     \tag{AE.11}
\end{equation}
Indeed FF.4 gives $K_0(Z)_{\mathbf Q}\simeq K_0(E)_{\mathbf Q}$
by restriction, using KST, levelwise nilpotent invariance and
Drinfeld's $K_{-1}(R)=0$ for henselian local rings
\citep[Theorem 3.7 and Remark (ii), p.~8]{frDrinfeld2004Tate}.
No normal-crossing hypothesis is used in that argument.
On $E$, $F_\gamma^{d+1}=0$ and $\psi^2$ acts by $2^j$ on
the $j$th associated graded piece.  Hence the projector onto
$S_d(E)$ is
\begin{equation}
       \prod_{j=0}^{d-1}
                \frac{\psi^2-2^j}{2^d-2^j}.           \tag{AE.12}
\end{equation}
For the regular algebraizable $Z$, rational Chern character
identifies this weight-$d$ subgroup with $\mathrm{CH}^d(Z)_{\mathbf Q}$
\citep[Chapter II, Proposition 4.9 and Corollary 8.9.1]
{D_Weibel2013KbookII}.  The finite bundles, cycles, relations
and operations descend to smooth finite-type pointed etale
models as in EP, and the isomorphism passes to their filtered
limit.  The same assertion applies to $Z\times P$.

Let $\theta_d$ be proper cycle pushforward
$\mathrm{CH}_1(E)_{\mathbf Q}\to\mathrm{CH}^d(Z)_{\mathbf Q}$
followed by $\rho_d$.  Then
\begin{align}
 \rho_d\bigl(\mathrm{CH}^d(Z)_{\mathrm{alg},\mathbf Q}\bigr)
       &=J_K^d(E),                                     \tag{AE.13}\\
 B^d(X)_{\mathbf Q}
       &\simeq S_d(E)\big/
           \bigl(J_K^d(E)+\theta_d\mathrm{CH}_1(E)_{\mathbf Q}\bigr).
                                                        \tag{AE.14}
\end{align}
Here is a proof of both inclusions, without a family-lifting
assumption.  Nilpotent invariance in characteristic $p$ makes
\[
 K_0(D\times P)_{\mathbf Q}\longrightarrow K_0(E\times P)_{\mathbf Q}
\]
an isomorphism of rational lambda rings
\citep[Theorem A]{rootGeisserHesselholt2011}; its inverse preserves
the rational gamma filtration, by the universal operations after
clearing denominators.  Lift an exceptional family in that
filtration and apply (AE.1).  Its two endpoint classes on $Z$
are forced by the point isomorphism to have pure weight $d$.
The actual cycle $\operatorname{ch}_d(b)$ therefore realizes
their difference as ordinary algebraic equivalence.  This proves
the inclusion of $J_K^d(E)$ in the left side of (AE.13).

Conversely a cycle family on $Z\times P$ has, on a smooth
finite-type model, an inverse rational Chern-character class
with that codimension-$d$ component and all others zero.  Its
pullback is an actual pure-weight-$d$ class in
$F_\gamma^dK_0(Z\times P)_{\mathbf Q}$.  Restricting it to
$E\times P$ gives the required exceptional perfect-family
relation.  This proves the opposite inclusion.  A smooth
nonproper parameter curve can be compactified and its finitely
many horizontal cycle components closed in the fixed product;
their torsion-free local coordinate rings over the regular
curve give flatness and preserve the two fibres.  Thus proper
curves generate exactly the ordinary relation group.
Finally the same closure in $Z\times P$ lifts relations from
$X$, and localization for Chow groups gives (AE.14).

This comparison removes the ambient endpoint-lifting problem in
every dimension in the displayed resolved scope.  It neither
asserts the existence of these resolutions in arbitrary positive
characteristic dimension nor computes $J_K^d(E)$ for a general
singular exceptional scheme.  In particular it does not replace
actual perfect families on $E\times P$ by independently moving
classes on its strata, and it does not assert that the full
negative-$K$ family obstruction vanishes.
